% ---- score2 standard preamble (from: python3 score2_tools.py preamble) ----
\documentclass[10pt]{extarticle}
\usepackage[a4paper,margin=1.5cm]{geometry}
\usepackage[T1]{fontenc}
\usepackage{amsmath,amssymb}
\usepackage{graphicx}
\usepackage[export]{adjustbox}
\usepackage{multicol}
\usepackage{enumitem}
\usepackage{booktabs,array,tabularx}
\setlength{\columnsep}{0.7cm}
\setlength{\parindent}{0pt}
\setlength{\parskip}{3pt}
\pagestyle{empty}
\makeatletter
% figure registry (figdefs.tex is generated by the extract step)
\newcommand{\qfigdef}[4]{\@namedef{qf@file@#1}{#2}\@namedef{qf@w@#1}{#3}\@namedef{qf@wide@#1}{#4}}
\newcommand{\qf@check}[1]{\@ifundefined{qf@file@#1}{\PackageError{score2}{Unknown figure id #1}{}}{}}
\newcommand{\qf@img}[4][]{\includegraphics[width=#2\dimexpr\@nameuse{qf@w@#4}pt\relax,max width=#3,max height=0.45\textheight,#1]{\@nameuse{qf@file@#4}}}
% Full-width questions: finalize wraps every question that contains a wide figure in \qfullbegin...\qfullend,
% so the WHOLE question (stem, figures, options) is set across the page in reading order, then the two columns
% resume. The question's head - its number, stem and first figure (\qfullsplit marks the end of the head) -
% is kept on one page, so a number is never left at a page foot with its figure overleaf; the rest may break
% across pages normally, so no big gaps are left. Without a \qfullsplit, the whole question is kept together
% if it is under 45% of a page. Never write these yourself.
\newif\ifqf@full
\newif\ifqf@boxed
\newsavebox{\qf@box}
\newcommand{\qf@open}{\global\qf@fulltrue\par\global\qf@boxedtrue
  \setbox\qf@box\vbox\bgroup\hsize\textwidth\linewidth\textwidth}
\newcommand{\qf@emit}[1]{\par\egroup\global\qf@boxedfalse
  \ifdim\dimexpr\ht\qf@box+\dp\qf@box\relax<#1\textheight\noindent\usebox{\qf@box}\par\else\unvbox\qf@box\fi}
\newcommand{\qfullsplit}{\ifqf@boxed\qf@emit{0.85}\fi}
\newcommand{\qf@close}{\ifqf@boxed\qf@emit{0.45}\fi\global\qf@fullfalse}
\newcommand{\qfullbegin}{\end{multicols}\qf@open}
\newcommand{\qfullend}{\qf@close\begin{multicols}{2}}
\newcommand{\qfullstart}{\qf@open}
\newcommand{\qfulldone}{\qf@close}
% \qfig[scale]{id}: block figure, centred, directly where it is written (size set by the tool).
% Fallback if a wide figure is not inside a full-width question: break out of the columns in place - no floats.
\newcommand{\qfig}[2][1]{\qf@check{#2}%
  \ifqf@full\qf@blk{#1}{#2}\else
  \ifnum\@nameuse{qf@wide@#2}=1 \expandafter\@firstoftwo\else\expandafter\@secondoftwo\fi
  {\end{multicols}\begin{center}\qf@img{#1}{\textwidth}{#2}\end{center}\begin{multicols}{2}}%
  {\qf@blk{#1}{#2}}\fi}
\newcommand{\qf@blk}[2]{\par\smallskip{\centering\qf@img{#1}{\linewidth}{#2}\par}\smallskip}
% \qfigi[scale]{id}: inline figure (inside option lists / tables / minipages); never breaks columns.
\newcommand{\qfigi}[2][1]{\qf@check{#2}\qf@img[valign=c]{#1}{\linewidth}{#2}}
\makeatother
% structure macros - lint reads these to check completeness and figure placement
\newcommand{\papersec}[1]{\par\bigskip{\centering\Large\bfseries PAPER-#1\par}\medskip}
\newcommand{\partsec}[2]{\par\medskip\noindent{\large\bfseries PART-#1: #2}\par\smallskip}
\newcommand{\sect}[1]{\par\smallskip\noindent\textbf{#1}\par}
\newcommand{\qn}[1]{\par\smallskip\noindent\textbf{#1.}~}
\newcommand{\paraq}[2]{\par\smallskip\noindent\textbf{Paragraph for Questions #1 and #2}\par}
\newenvironment{qoptions}{\begin{itemize}[leftmargin=1.8em,itemsep=1pt,topsep=2pt,label={}]}{\end{itemize}}
\newcommand{\opt}[2]{\item[(#1)] #2}
% figure registry (generated by score2_tools.py finalize)
\qfigdef{p1_phy_s1_q01_1}{TOAS - Test 04 (15 Mar 2022)-img/p1_phy_s1_q01_1.png}{53.4}{0}
\qfigdef{p1_phy_s1_q02_1}{TOAS - Test 04 (15 Mar 2022)-img/p1_phy_s1_q02_1.png}{118.8}{0}
\qfigdef{p1_phy_s1_q03_1}{TOAS - Test 04 (15 Mar 2022)-img/p1_phy_s1_q03_1.png}{50.6}{0}
\qfigdef{p1_phy_s2_q05_1}{TOAS - Test 04 (15 Mar 2022)-img/p1_phy_s2_q05_1.png}{120.2}{0}
\qfigdef{p1_phy_s2_q06_1}{TOAS - Test 04 (15 Mar 2022)-img/p1_phy_s2_q06_1.png}{74.7}{0}
\qfigdef{p1_phy_s2_q06_A}{TOAS - Test 04 (15 Mar 2022)-img/p1_phy_s2_q06_A.png}{80.6}{0}
\qfigdef{p1_phy_s2_q06_B}{TOAS - Test 04 (15 Mar 2022)-img/p1_phy_s2_q06_B.png}{87.0}{0}
\qfigdef{p1_phy_s2_q06_C}{TOAS - Test 04 (15 Mar 2022)-img/p1_phy_s2_q06_C.png}{84.5}{0}
\qfigdef{p1_phy_s2_q06_D}{TOAS - Test 04 (15 Mar 2022)-img/p1_phy_s2_q06_D.png}{85.0}{0}
\qfigdef{p1_phy_s2_q07_1}{TOAS - Test 04 (15 Mar 2022)-img/p1_phy_s2_q07_1.png}{132.7}{0}
\qfigdef{p1_phy_s2_q08_1}{TOAS - Test 04 (15 Mar 2022)-img/p1_phy_s2_q08_1.png}{148.5}{0}
\qfigdef{p1_phy_s2_q09_1}{TOAS - Test 04 (15 Mar 2022)-img/p1_phy_s2_q09_1.png}{287.3}{1}
\qfigdef{p1_phy_s2_q10_1}{TOAS - Test 04 (15 Mar 2022)-img/p1_phy_s2_q10_1.png}{126.8}{0}
\qfigdef{p1_phy_s3_para1_2_1}{TOAS - Test 04 (15 Mar 2022)-img/p1_phy_s3_para1_2_1.png}{263.0}{1}
\qfigdef{p1_phy_s3_para3_4_1}{TOAS - Test 04 (15 Mar 2022)-img/p1_phy_s3_para3_4_1.png}{40.1}{0}
\qfigdef{p1_phy_s3_para5_6_1}{TOAS - Test 04 (15 Mar 2022)-img/p1_phy_s3_para5_6_1.png}{165.6}{0}
\qfigdef{p1_phy_s4_q08_1}{TOAS - Test 04 (15 Mar 2022)-img/p1_phy_s4_q08_1.png}{101.8}{0}
\qfigdef{p1_chem_s1_q03_C}{TOAS - Test 04 (15 Mar 2022)-img/p1_chem_s1_q03_C.png}{57.5}{0}
\qfigdef{p1_chem_s1_q03_D}{TOAS - Test 04 (15 Mar 2022)-img/p1_chem_s1_q03_D.png}{51.7}{0}
\qfigdef{p1_chem_s2_q07_1}{TOAS - Test 04 (15 Mar 2022)-img/p1_chem_s2_q07_1.png}{166.5}{0}
\qfigdef{p1_chem_s2_q09_1}{TOAS - Test 04 (15 Mar 2022)-img/p1_chem_s2_q09_1.png}{209.1}{0}
\qfigdef{p1_chem_s2_q09_2}{TOAS - Test 04 (15 Mar 2022)-img/p1_chem_s2_q09_2.png}{88.0}{0}
\qfigdef{p1_chem_s2_q09_3}{TOAS - Test 04 (15 Mar 2022)-img/p1_chem_s2_q09_3.png}{88.2}{0}
\qfigdef{p1_chem_s2_q09_4}{TOAS - Test 04 (15 Mar 2022)-img/p1_chem_s2_q09_4.png}{85.8}{0}
\qfigdef{p1_chem_s2_q09_5}{TOAS - Test 04 (15 Mar 2022)-img/p1_chem_s2_q09_5.png}{84.2}{0}
\qfigdef{p1_chem_s2_q10_1}{TOAS - Test 04 (15 Mar 2022)-img/p1_chem_s2_q10_1.png}{199.6}{0}
\qfigdef{p1_chem_s3_para5_6_1}{TOAS - Test 04 (15 Mar 2022)-img/p1_chem_s3_para5_6_1.png}{89.9}{0}
\qfigdef{p1_chem_s4_q08_1}{TOAS - Test 04 (15 Mar 2022)-img/p1_chem_s4_q08_1.png}{224.4}{0}
\qfigdef{p2_phy_s1_q01_1}{TOAS - Test 04 (15 Mar 2022)-img/p2_phy_s1_q01_1.png}{111.1}{0}
\qfigdef{p2_phy_s1_q02_1}{TOAS - Test 04 (15 Mar 2022)-img/p2_phy_s1_q02_1.png}{147.9}{0}
\qfigdef{p2_phy_s1_q04_1}{TOAS - Test 04 (15 Mar 2022)-img/p2_phy_s1_q04_1.png}{120.1}{0}
\qfigdef{p2_phy_s1_q06_1}{TOAS - Test 04 (15 Mar 2022)-img/p2_phy_s1_q06_1.png}{156.5}{0}
\qfigdef{p2_phy_s2_para7_8_1}{TOAS - Test 04 (15 Mar 2022)-img/p2_phy_s2_para7_8_1.png}{107.5}{0}
\qfigdef{p2_phy_s2_para9_10_1}{TOAS - Test 04 (15 Mar 2022)-img/p2_phy_s2_para9_10_1.png}{72.3}{0}
\qfigdef{p2_phy_s3_para1_2_1}{TOAS - Test 04 (15 Mar 2022)-img/p2_phy_s3_para1_2_1.png}{95.5}{0}
\qfigdef{p2_phy_s3_para3_4_1}{TOAS - Test 04 (15 Mar 2022)-img/p2_phy_s3_para3_4_1.png}{36.2}{0}
\qfigdef{p2_phy_s3_para5_6_1}{TOAS - Test 04 (15 Mar 2022)-img/p2_phy_s3_para5_6_1.png}{118.0}{0}
\qfigdef{p2_phy_s4_q07_1}{TOAS - Test 04 (15 Mar 2022)-img/p2_phy_s4_q07_1.png}{76.9}{0}
\qfigdef{p2_phy_s4_q08_1}{TOAS - Test 04 (15 Mar 2022)-img/p2_phy_s4_q08_1.png}{103.0}{0}
\qfigdef{p2_phy_s4_q09_1}{TOAS - Test 04 (15 Mar 2022)-img/p2_phy_s4_q09_1.png}{158.4}{0}
\qfigdef{p2_chem_s1_q05_A}{TOAS - Test 04 (15 Mar 2022)-img/p2_chem_s1_q05_A.png}{186.8}{0}
\qfigdef{p2_chem_s1_q05_2}{TOAS - Test 04 (15 Mar 2022)-img/p2_chem_s1_q05_2.png}{132.5}{0}
\qfigdef{p2_chem_s1_q05_3}{TOAS - Test 04 (15 Mar 2022)-img/p2_chem_s1_q05_3.png}{101.4}{0}
\qfigdef{p2_chem_s1_q05_D}{TOAS - Test 04 (15 Mar 2022)-img/p2_chem_s1_q05_D.png}{207.4}{0}
\qfigdef{p2_chem_s1_q06_A}{TOAS - Test 04 (15 Mar 2022)-img/p2_chem_s1_q06_A.png}{153.9}{0}
\qfigdef{p2_chem_s1_q06_B}{TOAS - Test 04 (15 Mar 2022)-img/p2_chem_s1_q06_B.png}{78.2}{0}
\qfigdef{p2_chem_s1_q06_3}{TOAS - Test 04 (15 Mar 2022)-img/p2_chem_s1_q06_3.png}{81.8}{0}
\qfigdef{p2_chem_s1_q06_4}{TOAS - Test 04 (15 Mar 2022)-img/p2_chem_s1_q06_4.png}{27.5}{0}
\qfigdef{p2_chem_s1_q06_C}{TOAS - Test 04 (15 Mar 2022)-img/p2_chem_s1_q06_C.png}{63.4}{0}
\qfigdef{p2_chem_s1_q06_D}{TOAS - Test 04 (15 Mar 2022)-img/p2_chem_s1_q06_D.png}{143.3}{0}
\qfigdef{p2_chem_s2_para9_10_1}{TOAS - Test 04 (15 Mar 2022)-img/p2_chem_s2_para9_10_1.png}{128.3}{0}
\qfigdef{p2_chem_s2_q09_A}{TOAS - Test 04 (15 Mar 2022)-img/p2_chem_s2_q09_A.png}{86.2}{0}
\qfigdef{p2_chem_s2_q09_B}{TOAS - Test 04 (15 Mar 2022)-img/p2_chem_s2_q09_B.png}{69.4}{0}
\qfigdef{p2_chem_s2_q09_C}{TOAS - Test 04 (15 Mar 2022)-img/p2_chem_s2_q09_C.png}{66.8}{0}
\qfigdef{p2_chem_s2_q09_4}{TOAS - Test 04 (15 Mar 2022)-img/p2_chem_s2_q09_4.png}{206.5}{0}
\qfigdef{p2_chem_s3_q04_1}{TOAS - Test 04 (15 Mar 2022)-img/p2_chem_s3_q04_1.png}{183.9}{0}
\qfigdef{p2_chem_s4_q09_1}{TOAS - Test 04 (15 Mar 2022)-img/p2_chem_s4_q09_1.png}{49.2}{0}
\raggedcolumns
\makeatletter\@beginparpenalty=10000\makeatother
\newcolumntype{L}[1]{>{\raggedright\arraybackslash}p{#1}}
\begin{document}
\begin{center}{\Large\bfseries ALLEN TOAS --- Test 04 (15 Mar 2022)}\end{center}
\begin{multicols}{2}
\papersec{1}
\partsec{1}{PHYSICS}
\sect{SECTION-I(i)}
\qn{1} A cylinder is filled with a liquid of refractive index $\mu$. The radius of the cylinder is decreasing at a constant rate $K$. The volume of the liquid inside the container remains constant at $V$. The observer and the object $O$ are in a state of rest and at a distance $L$ from each other. The apparent velocity of the object as seen by the observer, is : (when radius of cylinder is $r$)
\qfig{p1_phy_s1_q01_1}
\begin{qoptions}
\opt{A}{$\dfrac{(1-\mu)2KV}{\pi\mu r^3}$}
\opt{B}{$\dfrac{(1-\mu)2KV}{\pi\mu Lr^2}$}
\opt{C}{$\dfrac{(1-\mu)2K}{\mu}$}
\opt{D}{$\dfrac{(1-\mu)K}{2\mu}$}
\end{qoptions}
\qn{2} A particle of mass 2 kg is moving along $x$-axis with a velocity of 5 m/s. It crosses the origin at $t=0$. A time varying force whose variation is as shown starts acting on a particle at $t=1$ sec. For this situation, mark out the correct statement.
\qfig{p1_phy_s1_q02_1}
\begin{qoptions}
\opt{A}{The particle will come to rest instantaneously at $t=2$ sec}
\opt{B}{The final velocity of the particle will be 18 m/s.}
\opt{C}{Work done by the force $F$ on the particle is 375 J.}
\opt{D}{All of the above}
\end{qoptions}
\qn{3} Two tennis balls of mass 60 g are attached with a massless rubber thread, and held in the vertical position as shown in the figure. In this position, the unstretched length of the rubber thread is 40 cm. The upper ball is slowly raised vertically upward, until the lower ball just becomes unsupported by the ground. At this time the length of the thread is 1 m. The rubber thread exerts a force which is proportional to its extension. Choose correct statement.
\qfig{p1_phy_s1_q03_1}
\begin{qoptions}
\opt{A}{Work done by external agent while the upper ball was raised is 0.6 J.}
\opt{B}{Work done by external agent while the upper ball was raised is 0.23 J.}
\opt{C}{Releasing the upper ball, the speed which it hits the lower one is 5.1 m/s.}
\opt{D}{The time that elapses between the release of the upper ball and the collision is 0.5 sec.}
\end{qoptions}
\qn{4} A mass of $m=1$ kg moves initially with a velocity of $v_0=10\sqrt{15}$ m/s. A constant power $P=1000$ Watt is applied so as to increase its velocity. Find the time that elapses before the acceleration is reduced to half of its original value.
\begin{qoptions}
\opt{A}{4 sec}
\opt{B}{9/4 sec}
\opt{C}{1 sec}
\opt{D}{4.5 sec}
\end{qoptions}
\sect{SECTION-I(ii)}
\qn{5} A small sphere of mass $m$ is connected by a string to a nail at $O$ and moves in a circle of radius $r$ on the smooth plane inclined at an angle $\theta$ with the horizontal. If the sphere has a velocity $u$ at the top position $A$. Mark the correct option(s) :-
\qfig{p1_phy_s2_q05_1}
\begin{qoptions}
\opt{A}{Minimum velocity at $A$ so that string does not get slack instantaneously is $\sqrt{\dfrac35 gr}$.}
\opt{B}{Tension at $B$ is $9mg/5$ if the sphere has required velocity as in option(A).}
\opt{C}{Tension at $C$ is $18\,mg/5$ in situation of option(B).}
\opt{D}{None of these}
\end{qoptions}
\qn{6} A semicircular wire of radius $R$ is oriented vertically. A small bead is released from rest from the top of the wire. It slides without friction under the influence of gravity to the bottom, where it then leaves the wire horizontally and falls distance $H$ to the ground. The bead lands a horizontal distance $D$ away from where it was launched. Which of the following is correct graph of $RH$ versus $D^2$ ?
\qfig{p1_phy_s2_q06_1}
\begin{qoptions}
\opt{A}{\qfigi{p1_phy_s2_q06_A}}
\opt{B}{\qfigi{p1_phy_s2_q06_B}}
\opt{C}{\qfigi{p1_phy_s2_q06_C}}
\opt{D}{\qfigi{p1_phy_s2_q06_D}}
\end{qoptions}
\qn{7} A block of mass 1 gm carrying a charge $1\,\mu$C is connected to a spring of stiffness $(\pi^2\times10^{-3})$ N/m whose other end is fixed to a wall. The block rests on a smooth ground as shown. By making a suitable arrangement, an electric field lying in the plane of the figure, directed towards right with strength $\left(\dfrac\pi{10}\right)$ MV/m is switched on exactly for an interval of 1 ms after every half second. Take the moment when the field is switched on for first time as $t=0$. Now, select the correct options.
\qfig{p1_phy_s2_q07_1}
\begin{qoptions}
\opt{A}{The maximum extension in the spring is 0.1414 m (nearly)}
\opt{B}{The spring never gets compressed.}
\opt{C}{The time taken by the block to start from mean position and reach the position of maximum elongation of spring is 0.75s.}
\opt{D}{The time taken by the block to start from mean position to the position of maximum elongation of spring is 0.5 s.}
\end{qoptions}
\qn{8} A point source of light $S$ is placed on the axis of the lens of focal length 20 cm as shown in the figure. A screen is placed normal to the axis of the lens at a distance $x$ from it. Treat all rays as paraxial. Then which of the following is/are correct regarding the intensity of the image produced on the screen of the point source ?
\qfig{p1_phy_s2_q08_1}
\begin{qoptions}
\opt{A}{As $x$ is increased from zero, intensity continuously decreases.}
\opt{B}{As $x$ is increased from zero, intensity first increases then decreases.}
\opt{C}{Intensity at centre of screen for $x=90$ cm and $x=110$ cm is same.}
\opt{D}{Radius of bright circle obtained on screen is equal to 2 cm for $x=300$ cm.}
\end{qoptions}
\qfullbegin
\qn{9} A concave mirror forms real image of a point object `$O$' lying on the optical axis at a distance of 60 cm from the pole of mirror. The focal length of mirror is 30 cm. Now, the mirror is cut into two halves symmetrically. In figure-1 both halves are shifted at a distance of 2 mm apart in a direction perpendicular to the optical axis. In figure-2 both halves are rotated (before being separated) about `$O$' with reference to original position. Let $I_1$ and $I_2$ are the image formed by both halves in both figure. Choose the correct statement(s).
\qfig{p1_phy_s2_q09_1}
\qfullsplit
\begin{qoptions}
\opt{A}{In figure-2, $I_1$ and $I_2$ are coincide.}
\opt{B}{In figure-1, $I_1$ and $I_2$ are separated.}
\opt{C}{In figure-1, the distance between $I_1$ and $I_2$ is 8 mm.}
\opt{D}{In figure-2, the distance between $I_1$ and $I_2$ is 4 mm.}
\end{qoptions}
\qfullend
\qn{10} Two plane mirrors are placed in such a way that they make an angle $50^\circ$ as shown in figure. A ray of light falls on one mirror and it reflects from each mirror once. For the given situation mark the correct statement(s).
\qfig{p1_phy_s2_q10_1}
\begin{qoptions}
\opt{A}{The value of $\theta$ must be less than $40^\circ$.}
\opt{B}{$\theta$ can have any value ranging from $0^\circ$ to $90^\circ$.}
\opt{C}{Total deviation produced in light ray is independent of value $\theta$.}
\opt{D}{Laws of reflection are valid for only plane and spherical mirrors.}
\end{qoptions}
\qfullbegin
\sect{SECTION-II(i)}
\paraq{1}{2}
An experiment is performed to calculate the focal length of a convex mirror using convex lens. The setup consists of two pins, optical bench, one convex lens and one convex mirror. When only convex lens is placed at a distance 60 cm from pin (object), its image is formed at a distance 30 cm from the lens as shown in figure-(i). Now convex mirror is also kept right to the lens without changing distance between pin (object) and lens. The mirror is kept at a distance 10 cm from lens such that position of image of object pin and position of object pin coincides as shown in figure-(ii)
\qfig{p1_phy_s3_para1_2_1}
\qfullsplit
Now answer the following questions.
\qfullend
\qn{1} Focal length of convex lens used (in cm) is :
\qn{2} Focal length of convex mirror (in cm) is :
\paraq{3}{4}
In the given figure, the block is attached with a system of three ideal springs $A$, $B$, $C$. The block is displaced slightly by a distance `$x$' from its equilibrium position vertically downwards and released.
\qfig{p1_phy_s3_para3_4_1}
\qn{3} The ratio of deformation of spring $B$ to spring $C$ is
\qn{4} The time period of small vertical oscillations of block is $2\pi\sqrt{a\dfrac mk}$. Find $a$
\paraq{5}{6}
A ray of light is incident normally on one face $30^\circ$--$60^\circ$--$90^\circ$ prism of refractive index $\mu_P=5/3$ immersed in water of refractive index $\mu_m=4/3$ as shown in the figure.
\qfig{p1_phy_s3_para5_6_1}
\qn{5} The exit angle $\theta_2$ of the ray is $\sin^{-1}\left(\dfrac x{16}\right)$, find the value of $x$.
\qn{6} The minimum refractive index of surrounding medium so that total internal reflection at point $P$ ceases is $\dfrac n{\sqrt3}$, find the value of $n$.
\sect{SECTION-II(ii)}
\qn{7} A biconvex thin lens is prepared from glass of refractive index $3/2$. The two bounding surfaces have equal radii of 35 cm each. One of the surfaces is silvered to make it reflecting. An object is placed infront of this lens in such a way that the image coincides with the object. If distance (in cm) of object from lens is $x$, then the value of $2x$ is
\qn{8} Two charged particles $(m,q)$ and $(2m,-2q)$ are placed in a gravity free space where a uniform electric field $E$ exists as shown in figure. After the particles are released, they stay at a constant distance from each other. What is the distance between them in centimeter? (Neglected the gravitational interaction). ($K$ is electrostatic constant and take $\dfrac{Kq}{E}=\dfrac{27}{8}$)
\qfig{p1_phy_s4_q08_1}
\qn{9} A circular plane mirror is kept on a horizontal surface, an opaque sphere of radius $R=20$ cm is placed symmetrically on the circular plane mirror. A small insect lies on the sphere at a vertical height of 36 cm from the circular mirror. Find the minimum radius (in cm) of mirror so that image of insect is formed by the mirror.
\partsec{2}{CHEMISTRY}
\sect{SECTION-I(i)}
\qn{1} In which of the following reaction bond cleavage of $\mathrm{B_2H_6}$ is different from the other reactions.
\begin{qoptions}
\opt{A}{$\mathrm{B_2H_6+2NMe_3\longrightarrow}$}
\opt{B}{$\mathrm{B_2H_6+2THF\longrightarrow}$}
\opt{C}{$\mathrm{B_2H_6+2CO\longrightarrow}$}
\opt{D}{$\mathrm{B_2H_6+2NH_3\longrightarrow}$}
\end{qoptions}
\qn{2} Initial rate data at $25^\circ$C are listed in the table for the reaction

$\mathrm{NH_{4(aq)}^{+}+NO_{2\,(aq)}^{-}\to N_{2(g)}+2H_2O_{(l)}}$
{\small\setlength{\tabcolsep}{3pt}\renewcommand{\arraystretch}{1.25}\begin{center}\begin{tabular}{|c|c|c|c|}
\hline
Expt.\ No. & Initial $\left[\mathrm{NH_4^+}\right]$ & Initial $\left[\mathrm{NO_2^-}\right]$ & \begin{tabular}[c]{@{}c@{}}Initial rate of\\ consumption of\\ $\mathrm{NH_4^+}$ M/s\end{tabular}\\
\hline
1. & 0.24 & 0.10 & $7.2\times10^{-6}$\\
\hline
2. & 0.12 & 0.10 & $3.6\times10^{-6}$\\
\hline
3. & 0.12 & 0.15 & $5.4\times10^{-6}$\\
\hline
\end{tabular}\end{center}}
What is the reaction rate when concentrations are $\left[\mathrm{NH_4^+}\right]=0.39$ M and $\left[\mathrm{NO_2^-}\right]=0.052$ M ?
\begin{qoptions}
\opt{A}{$6.1\times10^{-5}$ M/s}
\opt{B}{$6.1\times10^{-6}$ M/s}
\opt{C}{$6.1\times10^{-7}$ M/s}
\opt{D}{$3\times10^{-5}$ M/s}
\end{qoptions}
\qn{3} Which of the following is most reactive towards $\mathrm{S_N2}$ mechanism
\begin{qoptions}
\opt{A}{$\mathrm{CH_3{-}Cl}$}
\opt{B}{$\mathrm{CH_2{=}CH{-}CH_2{-}Cl}$}
\opt{C}{\qfigi{p1_chem_s1_q03_C}}
\opt{D}{\qfigi{p1_chem_s1_q03_D}}
\end{qoptions}
\qn{4} What is the final product of the given reaction

$\mathrm{CH_3{-}C{\equiv}CH}\xrightarrow[\mathrm{(ii)\ CH_3{-}CH_2{-}I}]{\mathrm{(i)\ Na}}\mathrm{Product}$
\begin{qoptions}
\opt{A}{$\mathrm{CH_2{=}CH{-}CH_2{-}CH_3}$}
\opt{B}{$\mathrm{CH_3{-}CH_2{-}C{\equiv}C{-}CH_2{-}CH_3}$}
\opt{C}{$\mathrm{CH{\equiv}C{-}CH_2{-}CH_2{-}CH_3}$}
\opt{D}{$\mathrm{CH_3{-}C{\equiv}C{-}CH_2{-}CH_3}$}
\end{qoptions}
\sect{SECTION-I(ii)}
\qn{5} Which of the following statement(s) is/are \textbf{CORRECT} -
\begin{qoptions}
\opt{A}{$\mathrm{[B_3O_6]^{-3}}$ ion is non-planar but $\mathrm{B_3N_3H_6}$ is planar}
\opt{B}{$\mathrm{(H_3Si)_3N}$ is planar but $\mathrm{(SiH_3)_3P}$ is pyramidal}
\opt{C}{$\mathrm{H_3C}$--NCS is bent but $\mathrm{SiH_3}$-NCS is linear}
\opt{D}{$\mathrm{(CH_3)_3N}$ is pyramidal but $\mathrm{(GeH_3)_3N}$ is planar}
\end{qoptions}
\qn{6} \begin{tabular}[t]{@{}L{3.3cm}@{\hspace{0.8em}}L{3.2cm}@{}}
Co-ordination entity & Wave length of light absorbed (nm)\\
$\mathrm{[CoCl(NH_3)_5]^{2+}}$ & P\\
$\mathrm{[Co(NH_3)_5(H_2O)]^{3+}}$ & Q\\
$\mathrm{[Co(NH_3)_6]^{3+}}$ & R\\
$\mathrm{[Co(CN)_6]^{3-}}$ & S
\end{tabular}

Select the correct orders.
\begin{qoptions}
\opt{A}{P $>$ R $>$ Q}
\opt{B}{P $>$ Q $>$ R $>$ S}
\opt{C}{P $>$ Q $>$ R}
\opt{D}{Q $>$ P $>$ R $>$ S}
\end{qoptions}
\qn{7} Equal mass of three gases He, $\mathrm{O_2}$ \& $\mathrm{SO_2}$ are present in a container initially as shown -
\qfig{p1_chem_s2_q07_1}
If mechanical stoppers are removed and movable pistons are allowed to attain equilibrium -

Choose the \textbf{incorrect} statement -
\begin{qoptions}
\opt{A}{Position of piston $A$ from the left end of the container is 15 cm}
\opt{B}{Position of piston $B$ from the left end of the container is 7 cm}
\opt{C}{Final pressure of $\mathrm{O_2}$ gas is 10 atm}
\opt{D}{Final pressure of He gas is 20 atm}
\end{qoptions}
\qn{8} For the following reaction

$\mathbf{2X+Y}\xrightarrow{\ k\ }\mathbf{P}$, the rate of reaction is $\dfrac{d[P]}{dt}=k[\mathbf X]$. Two moles of $\mathbf X$ are mixed with one mole of $\mathbf Y$ to make 1.0 L of solution. At 50 s, 0.5 mole of $\mathbf Y$ is left in the reaction mixture. The correct statement(s) about the reaction is(are) (Use: $\ln2=0.693$)
\begin{qoptions}
\opt{A}{The rate constant, $k$, of the reaction is $13.86\times10^{-4}\ \mathrm{s^{-1}}$.}
\opt{B}{Half--life of $\mathbf X$ is 50 s.}
\opt{C}{At 50 s, $-\dfrac{d[X]}{dt}=13.86\times10^{-3}\ \mathrm{mol\,L^{-1}\,s^{-1}}$.}
\opt{D}{At 100 s, $-\dfrac{d[Y]}{dt}=3.465\times10^{-3}\ \mathrm{mol\,L^{-1}\,s^{-1}}$.}
\end{qoptions}
\qn{9} \qfigi{p1_chem_s2_q09_1}

Correct pair of $X$ \& $Y$ respectively is :
\begin{qoptions}
\opt{A}{$X=\mathrm{PCl_5}$, $Y=$ \qfigi{p1_chem_s2_q09_2}}
\opt{B}{$X=\mathrm{SOCl_2}$, $Y=$ \qfigi{p1_chem_s2_q09_3}}
\opt{C}{$X=\mathrm{PCl_3}$, $Y=$ \qfigi{p1_chem_s2_q09_4}}
\opt{D}{$X=\mathrm{HCl/ZnCl_2}$, $Y=$ \qfigi{p1_chem_s2_q09_5}}
\end{qoptions}
\qn{10} \qfigi{p1_chem_s2_q10_1}
\begin{qoptions}
\opt{A}{Product contain $C_3$ axis of symmetry}
\opt{B}{Product is a bicyclo compound}
\opt{C}{Product contain plane of symmetry}
\opt{D}{Product can show stereoisomerism}
\end{qoptions}
\sect{SECTION-II(i)}
\paraq{1}{2}
Where two or more than two complexes having similar molecular weight (or) molecular formula but their structural arrangements are different these forms are known as isomer \& this phenomenon is known as isomerism; isomerism are of two types structural \& stereoisomerism ; stereoisomerism are of two types geometrical isomerism \& optical isomerism; compound which has plane of symmetry will not exhibit optical isomerism.
\qn{1} Total possible isomers of $\mathrm{K_4[Fe(CN)_6]}$.
\qn{2} Total geometrical isomers theoretically possible in an octatedral complex of formula $\mathrm{[Co(C_2O_4)(dmg)(en)]}$ (In this complex $\mathrm{dmg^{-1}}$ is bonded by both `N' atoms with central metal)
\paraq{3}{4}
For the first order dissociation reaction

$\mathrm{A_2B_3(aq.)\to2A^{3+}(aq.)+3B^{2-}(aq.)}$

a definite volume of solution is taken at different instant and its osmotic pressure rise is measured. Osmotic pressure rise at $t=0$ is 2 mm and at $t=10$ min is 6 mm. [Given density of solution is constant and equal to 1 gm/ml, $g=10\ \mathrm{m/s^2}$]
\qn{3} Calculate osmotic pressure of solution in mm at $t=20$ min.
\qn{4} Calculate time (in minute) for 87.5 \% of the reaction to complete -
\paraq{5}{6}
\qfig{p1_chem_s3_para5_6_1}
\qn{5} Value of `$n$' is?
\qn{6} Degree of unsaturation (D.O.U.) in product ($X$)?
\sect{SECTION-II(ii)}
\qn{7} Consider MOT for $\mathrm{B_2}$ \& $\mathrm{F_2}$ molecules \& find total number of electrons in both which is/are present in M.O. which is/are symmetrical around bond axis :-
\qn{8} A substance undergoes first order parallel reaction as shown. Calculate net activation energy (in kcal) if it is a known that $B$, $C$ and $D$ are formed in a molar ratio of $2:3:4$ respectively.
\qfig{p1_chem_s4_q08_1}
\qn{9} For a reaction $A\to P$

Half lite measured for two different values of initial concentrations $5\times10^{-3}$ M and $25\times10^{-4}$ M are 1 and 8 hrs respectively. If initial concentration is adjusted to $1.25\times10^{-3}$ M, what is the half-life? (in hrs)
\partsec{3}{MATHEMATICS}
\sect{SECTION-I(i)}
\qn{1} Value of $\displaystyle\lim_{x\to1}\dfrac{(\log(1+x)-\log2)(3\cdot4^{x-1}-3x)}{\left((7+x)^{1/3}-(1+3x)^{1/2}\right)\sin\pi x}$ is equal to
\begin{qoptions}
\opt{A}{$\dfrac9\pi(2\log2-1)$}
\opt{B}{$\dfrac9{4\pi}[\log4-1]$}
\opt{C}{$\dfrac9{2\pi}\left[\log4-\dfrac12\right]$}
\opt{D}{$\dfrac2{3\pi}[2\log2-1]$}
\end{qoptions}
\qn{2} If $\theta_1,\theta_2,\theta_3$ are three distinct real values such that \[\dfrac{\sin\theta_1+\sin\theta_2+\sin\theta_3}{\sin(\theta_1+\theta_2+\theta_3)}=\dfrac{\cos\theta_1+\cos\theta_2+\cos\theta_3}{\cos(\theta_1+\theta_2+\theta_3)}=2\] and $\cos(\theta_1+\theta_2)+\cos(\theta_2+\theta_3)+\cos(\theta_3+\theta_1)=\lambda$, then value of \[\lim_{x\to\lambda}\dfrac{\sqrt{x^2-\lambda^2}}{\sqrt{x-\lambda}+\sqrt{x-\sqrt\lambda}}\] is
\begin{qoptions}
\opt{A}{1}
\opt{B}{2}
\opt{C}{3}
\opt{D}{5}
\end{qoptions}
\qn{3} If $2f(\sin x)+\sqrt2f(\cos x)=\tan x$, then $\displaystyle\lim_{x\to1^-}(\sqrt{1-x})f((x))$ is - (Where $x\in\mathbb R-\left\{(2n+1)\dfrac\pi2;\ n\in I\right\}$)
\begin{qoptions}
\opt{A}{$\dfrac1{2\sqrt2}$}
\opt{B}{$\sqrt2$}
\opt{C}{$\dfrac1{\sqrt2}$}
\opt{D}{$\dfrac{\sqrt2}3$}
\end{qoptions}
\qn{4} The complete interval of values of $x$ in $\left(-\dfrac\pi2,\dfrac\pi2\right)$ satisfying the inequations $\cos x\cdot\cos2x>-\dfrac14$ and $3\cos x>1$, is
\begin{qoptions}
\opt{A}{$\left(0,\dfrac\pi3\right)$}
\opt{B}{$\left(-\dfrac\pi2,\dfrac\pi4\right)$}
\opt{C}{$\left(-\dfrac\pi3,\dfrac\pi3\right)$}
\opt{D}{$\left(-\dfrac\pi3,\dfrac{5\pi}{12}\right)$}
\end{qoptions}
\sect{SECTION-I(ii)}
\qn{5} If $\displaystyle\lim_{x\to0}\left(\dfrac{\tan x}x\right)^{\frac{\tan x}{x-\tan x}}$ exists and equal to $\ell$, then which of the following is (are) correct ?
\begin{qoptions}
\opt{A}{$\displaystyle\lim_{x\to\infty}\left(\dfrac{2x^2+1}{4x^2-1}\right)^{\frac{x^2}{2+x}}$ is less than $\ell$}
\opt{B}{If $f_1(x)=\dfrac x3+2\ \forall\ x\in\mathbb R$ and $f_n(x)$ is defined by $f_n(x)=f_1\left(f_{n-1}(x)\right)\ \forall\ n\geqslant2$ and $n\in\mathbb N$ then $\displaystyle\lim_{n\to\infty}f_n(x)$ is more than $\ell$}
\opt{C}{$\displaystyle\lim_{x\to\infty}\left(\dfrac{2x^2+1}{4x^2-1}\right)^{\frac{x^2}{2+x}}$ is more than $\ell$}
\opt{D}{If $f_1(x)=\dfrac x3+2\ \forall\ x\in\mathbb R$ and $f_n(x)$ is defined by $f_n(x)=f_1(f_{n-1}(x))\ \forall\ n\geq2$ and $n\in\mathbb N$ then $\displaystyle\lim_{n\to\infty}f_n(x)$ is less than $\ell$}
\end{qoptions}
\qn{6} Which of the following satisfy the equation $\sin^{-1}\sqrt{\dfrac{1+x}2}-\sqrt{2-x}=\cot^{-1}\left(\tan\sqrt{2-x}\right)-\sin^{-1}\sqrt{\dfrac{1-x}2}$
\begin{qoptions}
\opt{A}{$\left(2-\dfrac{\pi^2}4,1\right]$}
\opt{B}{$\left(2-\dfrac{\pi^2}4,0\right)$}
\opt{C}{$[-1,1]$}
\opt{D}{$\left[1-\dfrac{\pi^2}4,1\right]$}
\end{qoptions}
\qn{7} $f:\mathbb R\to\mathbb R$, $g:\mathbb R\to\mathbb R$, $h:\mathbb R\to\mathbb R$, defined such that $f\circ g\circ h:\mathbb R\to\mathbb R$ is a bijective function then :
\begin{qoptions}
\opt{A}{$h$ must be injective}
\opt{B}{$g$ may be many one into}
\opt{C}{$g$ must be surjective}
\opt{D}{$f$ must be surjective}
\end{qoptions}
\qn{8} A right angled triangle has one leg of length 1, another of length `$x$' \& hypotenuse of length `$y$'. The angle oppsoite to the side $x$ is $\theta$ then -
\begin{qoptions}
\opt{A}{$y^n-x^n=\dfrac{1-\sin^n\theta}{\cos^n\theta}$}
\opt{B}{$\displaystyle\lim_{\theta\to\frac\pi2}(y-x)=0$}
\opt{C}{$\displaystyle\lim_{\theta\to\frac\pi2}\left(y^2-x^2\right)=1$}
\opt{D}{$\displaystyle\lim_{\theta\to\frac\pi2}\left(y^3-x^3\right)$ does not exist.}
\end{qoptions}
\qn{9} Which of the following are incorrect ?
\begin{qoptions}
\opt{A}{Let $f$ be non constant periodic function with periods (need not be fundamental period) $T_1$ as well as $T_2$ then $\dfrac{T_1}{T_2}$ must be rational}
\opt{B}{Let $f$ and $g$ be periodic with periods $T_1$ and $T_2$ respectively such that LCM $(T_1,T_2)$ does not exist then $f+g$ must be non periodic}
\opt{C}{Let $f:\mathbb R\to\mathbb R$ be an invertible function then $f(x)=x\Leftrightarrow f(x)=f^{-1}(x)$}
\opt{D}{Let $f:\mathbb R\to\mathbb R$ be a function such that $f(x+y)=f(x)+f(y)\ \forall\ x,y\in\mathbb R$ then $f(x)$ must be $kx$ for some real number $k\in\mathbb R$.}
\end{qoptions}
\qn{10} The sum of infinite terms of series \[\begin{aligned}&\cot^{-1}\left(1^2+\sin^2\dfrac{2\pi}3\right)+\cot^{-1}\left(2^2+\sin^2\dfrac{4\pi}3\right)\\&\quad+\cot^{-1}\left(3^2+\sin^2\dfrac{8\pi}3\right)+\ldots\end{aligned}\] is
\begin{qoptions}
\opt{A}{$\cot^{-1}(1)$}
\opt{B}{$\tan^{-1}(2)$}
\opt{C}{$\dfrac12\sin^{-1}\left(\dfrac45\right)$}
\opt{D}{$\cot^{-1}(1)+\cot^{-1}(3)$}
\end{qoptions}
\sect{SECTION-II(i)}
\paraq{1}{2}
If $f(x)=4x^3-x^2-2x+1$ and \[g(x)=\left\{\begin{array}{ll}\min\{f(t);\ 0\leqslant t\leqslant x\} & ;\ \ 0\leqslant x\leqslant1\\[2pt] 3-x & ;\ \ 1<x\leqslant2\end{array}\right\},\] then
\qn{1} $g\left(\dfrac14\right)$ is equal to
\qn{2} $g\left(\dfrac34\right)$ is equal to
\paraq{3}{4}
Let $\alpha$ be a real root of $f(x)=8x^3+12x^2+6x+1$, then
\qn{3} The value of $\sin^{-1}(\sin[\alpha])+\cos^{-1}(\cos[\alpha])$ is equal to
\qn{4} The number of distinct real roots of equation $f(x)=0$, is ([.] denotes the greatest integer function)
\paraq{5}{6}
Let $a_1=1$ and $a_n=\sin(a_{n-1})$, $n>1$ \& $n\in\mathbb N$. If \[\lim_{n\to\infty}\dfrac{2^{2a_n}-2^{1+a_n}\cdot3^{a_n}+3^{2a_n}}{\cos a_n+1-e^{a_n}-e^{-a_n}}=-\dfrac\alpha3\,\ell n^{\beta}\left(\dfrac23\right),\] (where $\alpha$ and $\beta$ are prime numbers).
\qn{5} The value of $3\alpha$ is
\qn{6} The value of $2\beta$ is
\sect{SECTION-II(ii)}
\qn{7} If a sequence is such that $a_1=1$ and $a_n=n(1+a_{n-1})\ \forall\ n\geq2$ and $n\in\mathbb N$, then \[\lim_{n\to\infty}\ell n\left(\left(1+\dfrac1{a_1}\right)\left(1+\dfrac1{a_2}\right)\ldots\left(1+\dfrac1{a_n}\right)\right)\] is \underline{\hspace{2.5em}}.
\qn{8} If the solution set of $\sin^{-1}x+\sin^{-1}y\geq(x+y)^2-2[x(y+1)-(1-y)]+\pi$ is $(x_0,y_0)$ then $3x_0+y_0$ is
\qn{9} If $\displaystyle\lim_{x\to0}\dfrac{a+\log_e(1+x)+b\sin x}{x^2}=\ell$ (finite) then $|a+b+4\ell|$ (where $a,b,\ell\in\mathbb R$)
\newpage
\papersec{2}
\partsec{1}{PHYSICS}
\sect{SECTION-I(i)}
\qn{1} The refractive index of the medium with a certain region, $x>0$, $y>0$, changes with $y$. A thin light ray travelling in the $x$-direction in medium having refractive index $\mu_0=1$ strikes another medium of refractive index $\mu$ at right angles and moves through the medium along a circular arc of radius $R$ as shown in the figure. The material with the greatest known refractive index is diamond, but even the refractive index of this material does not reach the value $\mu_{\max}=2.5$. It is this limit that sets the maximum angular size of the arc the light ray can cover. Angular size of arc is the angle subtended by the arc at the centre.
\qfig{p2_phy_s1_q01_1}
\begin{qoptions}
\opt{A}{The variation of refractive index $\mu$ with $y$ is given as $\mu=\dfrac{R}{R-y}$}
\opt{B}{The unit vector in the direction of refracted light at $y=\dfrac R2$ is $\dfrac12\hat i+\dfrac{\sqrt3}{2}\hat j$}
\opt{C}{The maximum angular size of the arc of light is $\theta_{\max}$, the value of $\cos\theta_{\max}$ is $2/5$}
\opt{D}{The maximum angular size of the arc of light is $\theta_{\max}$, the value of $\cos\theta_{\max}$ is $3/5$}
\end{qoptions}
\qn{2} A particle of mass $m$ is placed on vertical fixed rough circular track and then it is given velocity $v$ vertically downward at position $A$ on track. If block moves on track with constant speed then
\qfig{p2_phy_s1_q02_1}
\begin{qoptions}
\opt{A}{Coefficient of friction between block and circular track as function of angle $\theta$ is $\mu=\dfrac{\sin\theta}{\cot\theta+\frac{v^2}{Rg}}$}
\opt{B}{Coefficient of friction between block and circular track as function of angle $\theta$ is $\mu=\dfrac{\sin\theta}{\cos\theta+\frac{v^2}{Rg}}$}
\opt{C}{Instantaneous power due to friction is $(-mg\sin\theta\,v)$}
\opt{D}{Work done from $A$ to $C$ by friction on block will be $(-mgR)$}
\end{qoptions}
\qn{3} A man is running with a constant acceleration (without slipping) on a plank which is placed on a horizontal smooth surface. Then choose the correct option :-
\begin{qoptions}
\opt{A}{Work done by friction on the man is negative.}
\opt{B}{Work done by friction on the man is positive.}
\opt{C}{Work done by friction on the plank is positive.}
\opt{D}{Work done by friction on the plank is negative.}
\end{qoptions}
\qn{4} The X-Y-Z space has an infinite charged sheet having surface charge density $\sigma$ on the Y-Z plane. A negative charge particle $-Q$ is placed at $(a,0,0)$. The variation of net electric field on $x$-axis is shown in the figure. Then choose the correct option(s) :
\qfig{p2_phy_s1_q04_1}
\begin{qoptions}
\opt{A}{$Q=\sigma\,2\pi a^2$}
\opt{B}{$Q=\sigma\,4\pi a^2$}
\opt{C}{The net electric field at the point $\left(\dfrac a2,0,0\right)$ is $\dfrac{5\sigma}{2\varepsilon_0}$}
\opt{D}{At point $\left(\dfrac{5a}2,0,0\right)$ the net electric field is $\dfrac{5\sigma}{18\varepsilon_0}$.}
\end{qoptions}
\qn{5} A point charge $q$ is kept at origin. The space around origin can be considered as non conducting sphere, tends to infinity and having continuous charge distribution, whose volume charge density is
\[\rho=\frac{-q}{4\pi a}\cdot\frac{e^{-r/a}}{r^2}\]
Where `$a$' is a positive constant and `$r$' is distance of a point from origin. The potential at infinity can be considered to be zero. The electric constant of space is 1. (Neglect the effect of polarization)
\begin{qoptions}
\opt{A}{The total charge of system is zero.}
\opt{B}{The total charge of system is $\dfrac q2$.}
\opt{C}{The magnitude of electrostatic field at a distance of $r$ from origin is $\dfrac{qe^{-r/a}}{4\pi\varepsilon_0r^2}$}
\opt{D}{The magnitude of electrostatic field at a distance of $r$ from origin is $\dfrac{qe^{-r/a}}{8\pi\varepsilon_0r^2}$}
\end{qoptions}
\qn{6} Consider the rays shown in the diagram as paraxial. The image of the virtual point object $O$ formed by the concave lens $LL'$ is
\qfig{p2_phy_s1_q06_1}
\begin{qoptions}
\opt{A}{Virtual.}
\opt{B}{Real.}
\opt{C}{Located below the principal axis.}
\opt{D}{Located left of the lens.}
\end{qoptions}
\sect{SECTION-I(ii)}
\paraq{7}{8}
A weightless rod of length $\ell$ with a small load of mass $m$ at one of its end is held vertical with its lower end hinged on a horizontal surface. The load touches a wedge of mass $M$ in this position. A slight jerk towards right sets the system in motion (see figure), with rod rotating freely in vertical plane about its lower end. There is no friction.
\qfig{p2_phy_s2_para7_8_1}
\qn{7} For what mass ratio $\dfrac Mm$ will the rod form an angle $\theta=\dfrac\pi3$ with the vertical at the moment the load separates from the wedge ?
\begin{qoptions}
\opt{A}{1}
\opt{B}{2}
\opt{C}{3}
\opt{D}{4}
\end{qoptions}
\qn{8} What is speed of the wedge at that moment ? Neglect friction.
\begin{qoptions}
\opt{A}{$\sqrt{\dfrac{g\ell}2}$}
\opt{B}{$\sqrt{\dfrac{g\ell}4}$}
\opt{C}{$\sqrt{\dfrac{g\ell}6}$}
\opt{D}{$\sqrt{\dfrac{g\ell}8}$}
\end{qoptions}
\paraq{9}{10}
A long helix made of thin wire is held vertical. The radius and pitch of the helix are $R$ and $\rho$ respectively. A bead begins to slide down the helix. Pitch is defined as displacement of particle along the axis of helix in one rotation. (Given : $\tan\theta=\dfrac{\rho}{2\pi R}$)
\qfig{p2_phy_s2_para9_10_1}
\qn{9} Find the normal force applied by the wire on the bead when the speed of the bead is $v$.
\begin{qoptions}
\opt{A}{$mg\cos\theta\sqrt{1+\left(\dfrac{v^2\cos\theta}{Rg}\right)^2}$}
\opt{B}{$mg\sin\theta\sqrt{1+\left(\dfrac{v^2\cos\theta}{Rg}\right)^2}$}
\opt{C}{$mg\sin\theta\sqrt{1+\left(\dfrac{v^2\sin\theta}{Rg}\right)^2}$}
\opt{D}{$mg\cos\theta\sqrt{1+\left(\dfrac{v^2\sin\theta}{Rg}\right)^2}$}
\end{qoptions}
\qn{10} Eventually, the bead acquires a constant speed of $v_0$. Find the coefficient of friction between the wire and the bead.
\begin{qoptions}
\opt{A}{$\mu=\dfrac{2\tan\theta}{\sqrt{1+\left(\frac{v^2\sin\theta}{Rg}\right)^2}}$}
\opt{B}{$\mu=\dfrac{\tan\theta}{\sqrt{1+\left(\frac{v^2\cos\theta}{Rg}\right)^2}}$}
\opt{C}{$\mu=\dfrac{\tan\theta}{2\sqrt{1+\left(\frac{v^2\cos\theta}{Rg}\right)^2}}$}
\opt{D}{$\mu=\dfrac{\tan\theta}{2\sqrt{1+\left(\frac{v^2\sin\theta}{Rg}\right)^2}}$}
\end{qoptions}
\sect{SECTION-II(i)}
\paraq{1}{2}
A tube of mass $M$ hangs from a thread and two balls of mass $m$ slide inside it without friction (see figure). The balls are released simultaneously from the top of the ring and slide down on opposite sides. $\theta$ defines the positions of balls at any time as shown in figure.
\qfig{p2_phy_s3_para1_2_1}
\qn{1} Minimum value of $m$ so that ring will start to rise is $kM$, find value of $k$.
\qn{2} If $M=0$, then the value of angle is $\theta$ at which the tube begins to rise. Find $\cos\theta$.
\paraq{3}{4}
A uniformly charged non conducting rod is suspended vertically at its end. The rod can swing freely in the vertical plane without any friction. The linear charge density on the rod is $\lambda$ C/m and it has a uniformly distributed mass of $M$. Length of the rod is $L$. A uniform horizontal electric field ($E$) is switched on in the region.
\qfig{p2_phy_s3_para3_4_1}
\qn{3} If value of electric field is $E_0$ so that the rod just manages to make itself horizontal. If $E_0=\dfrac{nMg}{\lambda L}$, find $n$.
\qn{4} If electric field $E_0$ is switched on, the maximum angular speed acquired by the rod during its motion is $\sqrt{\dfrac{xg}{L}}$. Find $x$.
\paraq{5}{6}
For a glass prism the plot of deviation ($\delta$) v/s angle of incidence ($i$) is as shown.
\qfig{p2_phy_s3_para5_6_1}
\qn{5} Find the refractive index of the glass.
\qn{6} Find the value of angle $i_1$ (in degree).
\sect{SECTION-II(ii)}
\qn{7} A heavy particle is attached to one end of a light string of length $\ell$ whose other end is fixed at $O$. The particle is projected horizontally with a velocity $v_0$ from its lowest position $A$. When the angular displacement of the string is more than $90^\circ$, the particle leaves the circular path at $B$. The string again becomes taut at $C$ such that $B,O,C$ are colinear. If $v_0^2=\left(\dfrac{4+n\sqrt2}{2}\right)g\ell$, find the value of $n$.
\qfig{p2_phy_s4_q07_1}
\qn{8} Three plane mirrors are kept as shown in the figure. A point object ($O$) is kept at the centroid of the triangle seen in the figure. How many images will be formed ?
\qfig{p2_phy_s4_q08_1}
\qn{9} A pencil ($AB$) of length 20 cm is moving along the principal axis of a concave mirror $MM'$, with a velocity 5 m/s approaching the mirror. The mirror itself is moving away from the pencil at a speed of 2 m/s. Radius of curvature of mirror is 40 cm. If the rate of change of length of the image of the pencil at the instant end $A$ is at a distance of 60 cm from the mirror is $\dfrac x4$ m/s, find $x$.
\qfig{p2_phy_s4_q09_1}
\partsec{2}{CHEMISTRY}
\sect{SECTION-I(i)}
\qn{1} Which complex of the following pairs has the larger value of $\Delta_0$ ?\\
(i) $\mathrm{[Co(CN)_6]^{3-}}$ and $\mathrm{[Co(NH_3)_6]^{3+}}$\\
(ii) $\mathrm{[Co(H_2O)_6]^{3+}}$ and $\mathrm{[Rh(H_2O)_6]^{3+}}$\\
(iii) $\mathrm{[Co(H_2O)_6]^{3+}}$ and $\mathrm{[Co(H_2O)_3]^{3+}}$\\
(iv) $\mathrm{[Co(NH_3)_6]^{3+}}$ and $\mathrm{[CoF_6]^{3-}}$\\
Select the correct one
\begin{qoptions}
\opt{A}{$\mathrm{[Co(CN)_6]^{3-}>[Co(H_2O)_6]^{3+}}$}
\opt{B}{$\mathrm{[Co(H_2O)_6]^{2+}<[Co(H_2O)_6]^{3+}}$}
\opt{C}{$\mathrm{[Co(H_2O)_6]^{3+}>[Rh(H_2O)_6]^{3+}}$}
\opt{D}{$\mathrm{[Co(NH_3)_6]^{3+}<[CoF_6]^{3-}}$}
\end{qoptions}
\qn{2} Which of the following given pair of species exhibit similar kind of magnetic behaviour as well as at least one same kind of isomerism ?
\begin{qoptions}
\opt{A}{$\mathrm{[Zn(NH_3)_4]Cl_2}$ and $\mathrm{[Co(H_2O)_6]Cl_2}$}
\opt{B}{$\mathrm{[Cr(H_2O)_4Cl_2]^{+1}}$ and $\mathrm{[Co(NH_3)_3Cl_3]^{-1}}$}
\opt{C}{$\mathrm{K_3[Co(C_2O_4)_3]}$ and $\mathrm{K[Fe(EDTA)]}$ ; (EDTA = $\mathrm{EDTA^{4-}}$)}
\opt{D}{$\mathrm{[Co(NH_3)_5(NO_2)]SO_4}$ and $\mathrm{[Pd(PPh_3)_2(SCN)_2]}$}
\end{qoptions}
\qn{3} In an evacuated vessel of 8.2 L capacity, 0.2 moles of $\mathrm{A(g)}$ is taken, vessel is sealed and heated to $127^\circ$C, by which the following first order reaction occurred at constant volume and temperature :-\\
$\mathrm{A(g)}\rightarrow\mathrm{2B(g)+C(s)}$\\
The vapour pressure of $\mathrm{C(s)}$ is 0.1 atm at $127^\circ$C. If the half life of reaction is 20 min, then the correct information(s) related with reaction is/are : ($R=0.082$ L-atm/K-mol)
\begin{qoptions}
\opt{A}{$P_{\mathrm{total}}=1.7$ atm\quad at $t=\infty$}
\opt{B}{$P_{\mathrm{total}}=0.98$ atm after 10\% decomposition of $\mathrm{A(g)}$}
\opt{C}{$P_{\mathrm{total}}=0.88$ atm after 5\% decompositon of $\mathrm{A(g)}$}
\opt{D}{$P_{\mathrm{total}}=1.3$ atm at $t=20$ min}
\end{qoptions}
\qn{4} 1.22 gm $\mathrm{C_6H_5COOH}$ is added into 100 gm of acetone, elevation in B.P. is found to be $0.17\,^\circ$C. If $K_b$ (acetone) = 1.7 K - Kg mol$^{-1}$ then select the incorrect options
\begin{qoptions}
\opt{A}{Observed molar mass of benzoic acid is 61 gm/mole.}
\opt{B}{Observed molar mass of benzoic acid is 122 gm/mole}
\opt{C}{Benzoic acid neither dissociate nor associate in acetone}
\opt{D}{Benzoic acid exist as dimer in acetone due to intermolecular H-bonding}
\end{qoptions}
\qn{5} Identify reactions correctly matched with their major products ?
\begin{qoptions}
\opt{A}{\qfigi{p2_chem_s1_q05_A}}
\opt{B}{$\mathrm{Me{-}C{\equiv}CH}$ \qfigi{p2_chem_s1_q05_2}}
\opt{C}{$\mathrm{Me{-}C{\equiv}CH}$ \qfigi{p2_chem_s1_q05_3}}
\opt{D}{\qfigi{p2_chem_s1_q05_D}}
\end{qoptions}
\qn{6} Identify the reaction(s) with correct formation of major product ?
\begin{qoptions}
\opt{A}{\qfigi{p2_chem_s1_q06_A}}
\opt{B}{\qfigi{p2_chem_s1_q06_B}\qfigi{p2_chem_s1_q06_3}}
\opt{C}{\qfigi{p2_chem_s1_q06_4}$\xrightarrow[\mathrm{(2)\ H_2O}]{\mathrm{(1)\ CH_3{-}MgBr}}$\qfigi{p2_chem_s1_q06_C}}
\opt{D}{\qfigi{p2_chem_s1_q06_D}}
\end{qoptions}
\sect{SECTION-I(ii)}
\paraq{7}{8}
2 moles of a volatile liquid (A) and 3 moles of another volatile liquid (B) are present in a solution at $25^\circ$C temperature in a closed container. If vapour pressure of pure A and pure B at $25^\circ$C are $P^0_A=300$ torr, $P^0_B=200$ torr respectively (above solution follows Raoult's law of partial vapour pressure)
\qn{7} Calculate vapour pressure of distillate of the solution at $25^\circ$C temperature.
\begin{qoptions}
\opt{A}{250 torr}
\opt{B}{240 torr}
\opt{C}{280 torr}
\opt{D}{270 torr}
\end{qoptions}
\qn{8} Which of the following is incorrect ?
\begin{qoptions}
\opt{A}{Above solution form an Azeotropic mixture}
\opt{B}{Above solution can be separated by fractional distillation}
\opt{C}{Mole fraction of A in the 1st drop of vapour at $25^\circ$C is 0.5}
\opt{D}{Vapour pressure of solution is 240 torr at $25^\circ$C}
\end{qoptions}
\paraq{9}{10}
Alkyl halides can undergo $\mathrm{S_N^{\,1}}$, $\mathrm{S_N^{\,2}}$, $\mathrm{E_1}$, $\mathrm{E_2}$, $\mathrm{E_{1cb}}$ depending on various conditions.
\qfig{p2_chem_s2_para9_10_1}
\qn{9} Identify `Y' :
\begin{qoptions}
\opt{A}{\qfigi{p2_chem_s2_q09_A}}
\opt{B}{\qfigi{p2_chem_s2_q09_B}}
\opt{C}{\qfigi{p2_chem_s2_q09_C}}
\opt{D}{Both (B) \& (C)}
\end{qoptions}
\qn{10} \qfig{p2_chem_s2_q09_4}
Select the correct statement :
\begin{qoptions}
\opt{A}{Formation of U follows $\mathrm{E^2}$ mechanism to give saytzeff product}
\opt{B}{Formation of U follows $\mathrm{E_{1cb}}$ mechanism to give saytzeff product}
\opt{C}{Formation of U follows $\mathrm{E^1}$ mechanism to give Hoffmann product}
\opt{D}{Formation of U follows $\mathrm{E^2}$ mechanism to give Hoffmann product}
\end{qoptions}
\sect{SECTION-II(i)}
\paraq{1}{2}
The metal--carbon bond in metal carbonyls posses both $\sigma$ \& $\pi$ character. The metal to ligand bonding creates a synergic effect which strengthens the bond between CO and the metal. Due to synergic effect M--C bond length \& C--O bond length affected.
\qn{1} Consider $\mathrm{[Co(CO)_4]^-}$ as $x$ and $\mathrm{[Mn(CO)_6]^+}$ as $y$ then find how many statements are correct?\\
(i) The C -- O bond length in $x$ is more than that in $y$\\
(ii) The C -- O bond energy in $x$ is more than that in $y$\\
(iii) Double bond character in metal to carbon bond in $x$ has more than that in $y$\\
(iv) Both $x$ and $y$ have the same EAN of the central metal atom
\qn{2} Consider the following species\\
$\mathrm{[Ni(CO)_4]}$, $\mathrm{[V(CO)_6]^-}$, $\mathrm{CO^+}$, $\mathrm{CO_2}$, $\mathrm{H_2CO_3}$, $\mathrm{CH_3COO^-}$, $\mathrm{[NiP(CH_3)_3(CO)_3]}$, $\mathrm{[Cr(CO)_6]}$, $\mathrm{[NiPF_3(CO)_3]}$\\
If bond length in CO is `$d$' in isolated molecule.\\
The number of species having higher C-O bond length then `$d$' is/are `$X$'\\
The number of species having lower C-O bond length then `$d$' is/are `$Y$'\\
Find the value of $\left(\dfrac XY\right)$
\paraq{3}{4}
A hypothetical metal $M$ (Atomic weight = 450) crystallizes as a cubic crystal such that metal atoms touch each other along body diagonal of cube. Edge length of cubic unit cell is 1.0 nm. If this metal is melted then it's density in molten form is 1.2 gm/cc. ($N_A=6\times10^{23}$)
\qn{3} Density of solid metal in gm/cc is -
\qn{4} Percentage of empty space in molten state of metal will be -
\paraq{5}{6}
\qfig{p2_chem_s3_q04_1}
\qn{5} D.O.U. of compound B $=n_1$ and D.O.U. of compound C $=n_2$ then value of $|n_1-n_2|$
\qn{6} How many option(s) is/are correct\\
(i) (C) gives turbidity in 5 min with lucas reagent\\
(ii) (C) gives racemic mixture in product with $\mathrm{PCl_3}$\\
(iii) (C) is isomers of anisole\\
(iv) (B) gives $\mathrm{CO_2}$ gas with $\mathrm{NaHCO_3}$
\sect{SECTION-II(ii)}
\qn{7} If $x$ and $y$ are total number of electrons which are present in non-axial and axial set of d-orbitals respectively on central metal ion of $\mathrm{[Ni(DMG)_2]}$, then calculate value of $\dfrac{2x^2}{y}$ \underline{\hspace{1.5cm}}.
\qn{8} $\alpha$-particle $\left({}^{4}_{2}\mathrm{He}^{++}\right)$ is moving along $x$-axis with 0.4 meV. If uncertainty in velocity is 1 \% and minimum uncertainty in x-co-ordinate is $1.80\times10^{-x}$ m, then $x$ is -
\qn{9} Consider the following complex for given three statements :\\
dichlorobis(ethylenediamine)chromium(III) chloride\\
(i) Number of unpaired electrons in complex $=(X)$\\
(ii) One mole of complex \qfigi{p2_chem_s4_q09_1} $(Y)$ mole of AgCl.\\
(iii) Total stereo isomeric forms of above complex $=(Z)$ Find $X+Y+Z$
\partsec{3}{MATHEMATICS}
\sect{SECTION-I(i)}
\qn{1} Let $(1+\sqrt2)^n=A_n+B_n\sqrt2$, where $A_n$ and $B_n$ are rational numbers and $n\in N$, then the value of $\displaystyle\lim_{n\to\infty}\dfrac{A_n}{B_n}$ is equal to -\\
(where $[.]$ denotes greatest integer function)
\begin{qoptions}
\opt{A}{$\left[\displaystyle\lim_{x\to0}\dfrac{\sin x}{x}\right]$}
\opt{B}{$\displaystyle\lim_{x\to0}\left[\dfrac{\sin x}{x}\right]$}
\opt{C}{$\displaystyle\lim_{x\to\pi/4}\dfrac{|\cos x-\sin x|}{\left|x-\frac\pi4\right|}$}
\opt{D}{$\left[\displaystyle\lim_{x\to0}\dfrac{\sin\sqrt2\,x}{x}\right]$}
\end{qoptions}
\qn{2} $\displaystyle\lim_{x\to-\infty}\left(\sqrt{x^2-\dfrac{4x}{3}}+x+2\right)$ is equal to-
\begin{qoptions}
\opt{A}{$\dfrac43$}
\opt{B}{$-\dfrac43$}
\opt{C}{$\dfrac83$}
\opt{D}{$-\dfrac83$}
\end{qoptions}
\qn{3} Given $\displaystyle\lim_{x\to\alpha}\dfrac{\sin\left(\left[x^3\right]-[x]^3\right)}{(x-\alpha)^3}$ exists finitely, then the value of $\alpha$ can be (where $\alpha\in I$ \& $[.]$ denotes greatest integer function)-
\begin{qoptions}
\opt{A}{0}
\opt{B}{1}
\opt{C}{$-1$}
\opt{D}{2}
\end{qoptions}
\qn{4} Let \[f(x)=\left\{\begin{array}{ll}\min\left\{\sin^{-1}x,\cos^{-1}x\right\}&;\ 0\le x\le1\\ \max.\left\{\sin^{-1}x,\cos^{-1}x\right\}&;\ -1\le x<0\end{array}\right.\] and $f(x)=\dfrac12$, $f(x)=\dfrac14$ \& $f(x)=2$ has number of solution(s) $\lambda_1,\lambda_2,\lambda_3$ respectively, then
\begin{qoptions}
\opt{A}{$\lambda_1=2$}
\opt{B}{$\lambda_2=1$}
\opt{C}{$\lambda_2=2$}
\opt{D}{$\lambda_3=1$}
\end{qoptions}
\qn{5} Let $f:R\to R$ defined by $f(x)=\cos^{-1}(-\{-x\})$, where $\{x\}$ is fractional part function, then which of the following is/are true ?
\begin{qoptions}
\opt{A}{$f(x)$ is many one function}
\opt{B}{period of $f(x)$ is 1}
\opt{C}{graph of $f(x)$ lies above $x$-axis}
\opt{D}{$f(x)$ is neither even nor odd function}
\end{qoptions}
\qn{6} If a function $f:R\to R$ satisfies $\left|\displaystyle\sum_{k=1}^{n}3^k\left(f(x+ky)-f(x-ky)\right)\right|\le1$ for every $n\in N$ and for all $x,y\in R$, then which of the following is/are true
\begin{qoptions}
\opt{A}{$\left|3^n\left(f(x+ny)-f(x-ny)\right)\right|\le2$}
\opt{B}{$|f(u)-f(v)|\le\dfrac{2}{3^n}$}
\opt{C}{$f(x)$ is a constant function}
\opt{D}{$f(x)$ is an odd function}
\end{qoptions}
\sect{SECTION-I(ii)}
\paraq{7}{8}
If $\tan^2\alpha\tan^2\beta+\tan^2\beta\tan^2\gamma+\tan^2\gamma\tan^2\alpha+2\tan^2\alpha\tan^2\beta\tan^2\gamma=1$.\\
(where $\alpha,\beta,\gamma\in R-\left\{\dfrac{n\pi}2\right\}$, $n\in I$)
\qn{7} Value of $\left(\Pi\cot^2\alpha-\Sigma\cot^2\alpha\right)$ is equal to -
\begin{qoptions}
\opt{A}{0}
\opt{B}{$-1$}
\opt{C}{1}
\opt{D}{2}
\end{qoptions}
\qn{8} Value of $\cos2\alpha+\cos2\beta+\cos2\gamma$ is equal to -
\begin{qoptions}
\opt{A}{$-1$}
\opt{B}{0}
\opt{C}{1}
\opt{D}{2}
\end{qoptions}
\paraq{9}{10}
The sequence $\{a_n\}_{n=1}^{+\infty}$ is defined by $a_1=0$ and $a_{n+1}=a_n+4n+3$, $n\ge1$
\qn{9} The remainder when $a_{50}$ is divided by 50, is
\begin{qoptions}
\opt{A}{0}
\opt{B}{3}
\opt{C}{29}
\opt{D}{47}
\end{qoptions}
\qn{10} The value of
\[\lim_{n\to\infty}\frac{\sqrt{a_n}+\sqrt{a_{4n}}+\sqrt{a_{4^2n}}+\sqrt{a_{4^3n}}+\cdots+\sqrt{a_{4^{10}n}}}{\sqrt{a_n}+\sqrt{a_{2n}}+\sqrt{a_{2^2n}}+\sqrt{a_{2^3n}}+\cdots+\sqrt{a_{2^{10}n}}}\]
equals
\begin{qoptions}
\opt{A}{685}
\opt{B}{683}
\opt{C}{684}
\opt{D}{680}
\end{qoptions}
\columnbreak
\sect{SECTION-II(i)}
\paraq{1}{2}
Consider $f(x)=\tan^{-1}\left(\dfrac{|x|}{\sqrt{1-x^2}}\right)+\cos^{-1}\left(\dfrac1{\sqrt{1+x^2}}\right)$ then
\qn{1} Number of integers in the range of $f(x)$ is
\qn{2} For $\theta\in\left(\dfrac\pi4,\dfrac\pi2\right)$, number of solution satisfying $f(\sin\theta)<f(\cos\theta)$ is
\paraq{3}{4}
Let $a=\sin10^\circ$, $b=\sin50^\circ$, $c=\sin70^\circ$.
\qn{3} Value of \[abc\left(\dfrac{a+b}{c}\right)\left(\dfrac1a+\dfrac1b-\dfrac1c\right)\] is equal to -
\qn{4} Value of $\dfrac{a^4+b^4+c^4}{(ab)^2+(bc)^2+(ca)^2}$ is -
\paraq{5}{6}
Consider the function $f:A\to[0,5]$ $f(x)=\sqrt{9-(x-1)(x-3)(x-5)(x-7)}$.\\
(Where set $A$ represents all values of $x$ for which $f(x)$ is defined.)
\qn{5} In the domain of the function, number of possible integral values of $x$ is equal to
\qn{6} Among the following, the number of correct statements is\\
(i) function is many one \quad (ii) function is one-one\\
(iii) function is onto \quad (iv) function is into
\sect{SECTION-II(ii)}
\qn{7} Let $f(x)$ be a function satisfying $f(x)+f(x-1)=x^2$, $\forall\,x\in R$. If $f(19)=94$ then the remainder when $f(94)$ is divided by 1000 is
\qn{8} Let $f(x)=a_0+a_1x+a_2x^2+a_3x^3+a_4x^4$ where $a_i\in R$ and $\displaystyle\lim_{x\to0}\left(\dfrac{f(x)-x+2x^2}{3x^3}\right)^{2/x}=e^4$, then $\displaystyle\sum_{i=0}^{4}a_i$ is equal to
\qn{9} A function $f$ satisfies $f(3x)=3f(x)$ $\forall\,x\in[1,\infty)$ and $f(x)=1-|x-2|$ for $x\in[1,3]$. The smallest $x$ for which $f(x)=f(2001)$ is
\end{multicols}
\end{document}
